\scanpage{3}
\mspage{2}
6. \uline{Niveau d'un motif} \inserted{Résulte de la présence de $M^+$ dans $M$.} Un motif $M$ est dit \uline{de niveau} $\leq i$ s'il est de la forme $M'(n)$, où $M'\in\operatorname{Ob}M^+(k)$ est à composantes homogènes de degrés $\leq i$.
\margin{Exemple : si $i>n=\dim X$, alors $h^i(X)$ est de \uline{niveau} $\leq 2n-i<i$}
On dit que $M$ est purement de niveau $i$ s'il est de niveau $\leq i$, et si aucun sous-motif $\ne0$ n'est de niveau $\leq i-1$. \inserted{Cet $i$ est unique et $\geq0$.} Tout motif \inserted{$M$} se décompose \inserted{de façon unique} en somme de motifs purement de niveau $i$
\[
M=\coprod_{i\in\mathrm Z}M_i
\]
donc s'il est de \deleted{degré} poids $n$, \deleted{il s'écrit} \inserted{les $i$ \deleted{\unsure{tels}} que $M_i\ne0$ ont la $=$ parité que $n$, et}
\[
M=\coprod_i M'_i\left(\frac{\unsure{n}-i}{2}\right)\qquad M'_i\in\mathord{\tikz[baseline=(c.base)]{\node[draw,ellipse,inner sep=2pt](c){$M^{+i}_{\mathrm{pur}}(k)$};}}
\]
où $M'_i$ est \deleted{\unsure{purement de poids}} \inserted{\uline{effectif}} pur de poids et niveau $i$.

\deleted{\unsure{Il est facile}} Les motifs de \deleted{poids}\inserted{niveau} $\leq i$ sont ceux qui, modulo twist de Tate, ``proviennent'' de variétés $X\in V(k)$ de dimension $\leq i$. Ceci résulte du fait que $h(X)$ est de niveau $\leq\dim X$, et que \deleted{\unsure{si $i$|Addition}} si $i\leq\dim X$, $h^i(X)\unsure{\longleftrightarrow|\lhook\joinrel\longrightarrow}h^i(Y)$, où $Y$ est une section linéaire de dim $i$ (\unsure{en vertu|au sens|sens} du ``th. de Lefschetz'').
\margin{\uline{Ex} Si $i\geq n=\dim X$ $h^i(X)$ est de niveau $\leq2n-i$}
Autre interprétation en termes de cohomologie $l$-adique \deleted{(en supposant $k$ alg. clos, ce qui est loisible, car la notion de niveau, pureté du \unsure{m} niveau, est invariante par extension du corps de base).} On regarde dans $H^i(\bar X,Q_\ell)$ la ``filtration géométrique'' définie par
\[
\mathrm{Filt}^{\unsure{j|J}}H^i(\bar X,Q_\ell)=
\underset{U,\;U\text{ ouvert tel que }\operatorname{codim}(X-U)\geq \unsure{j|d}}{\underline{\lim}}\,
\operatorname{Ker}\bigl(H^i(\bar X,Q_\ell)\to H^i(U,Q_\ell)\bigr)
\]
% archive sheet number: 2

